\begin{statementbox}
	\textbf{\textcolor{CStatement}{条件}}
	\begin{description}[style=nextline, leftmargin=2.8em, labelsep=0.5em, itemsep=0.35em]
		\item[\textbf{\textcolor{CStatement}{条件 1（集合与元素）：}}]
			设 $A$ 为任意集合，$a,b$ 为任意元素。
	\end{description}
	\textbf{\textcolor{CStatement}{结论}}
	\begin{description}[style=nextline, leftmargin=2.8em, labelsep=0.5em, itemsep=0.45em]
		\item[\textbf{\textcolor{CStatement}{结论一（确定性定义）：}}]
			\textbf{\textcolor{CStatement}{直观描述：}}
			任意对象要么属于该集合，要么不属于，绝不含糊。符号形式：$\forall x$，$x\in A$ 与 $x\notin A$ 有且仅有一个成立。
		\item[\textbf{\textcolor{CStatement}{结论二（互异性定义）：}}]
			\textbf{\textcolor{CStatement}{直观描述：}}
			集合中重复出现的元素视为同一个，不额外计数。
			\[
				\{a,a\} = \{a\}.
			\]
		\item[\textbf{\textcolor{CStatement}{结论三（无序性定义）：}}]
			\textbf{\textcolor{CStatement}{直观描述：}}
			元素书写的先后顺序不影响集合本身的含义。
			\[
				\{a,b\} = \{b,a\}.
			\]
		\item[\textbf{\textcolor{CStatement}{常见变形（集合书写规范）：}}]
			\textbf{\textcolor{CStatement}{直观描述：}}
			列举元素时须遵守互异且顺序无关。例如，
			\[
				\{1,2\} = \{2,1\},\quad \{1,1,2\} = \{1,2\}.
			\]
	\end{description}
	\textbf{\textcolor{CStatement}{使用提醒：}}
	三大特性是判断集合相等、子集关系以及进行集合运算的前提，忽视任意特性均会导致概念错误。
\end{statementbox}
